\chapter{Unified Multi-Robot Access and Control}
\label{chap:multi-robot}

\section{What ``unified'' means}

In this project, a \emph{unified} access layer is defined by four concrete
properties:

\begin{enumerate}
  \item common YAML configuration declares selectable robot instances;
  \item a common launch path spawns the selected instances and their bridges;
  \item a separate common command helper validates a named robot and its joint
  vector;
  \item instance-derived namespaces isolate commands and feedback.
\end{enumerate}

The common layer preserves the differences that matter to operation. Fixed-base
arms use joint trajectories, while TIAGo additionally exposes base velocity and a
linear torso joint. Joint limits, counts, names and units remain
robot-specific. The common joint-command helper validates the selector, vector
and units for \path{/joint_trajectory}; interfaces specific to a particular
robot structure, such as \path{/cmd_vel}, remain separate.

\section{Multi-robot operation in one Gazebo process}

The first phase established per-instance command interfaces for heterogeneous robots within
one Gazebo process. Repeatable operation depends on a consistent identity
chain. Each selected YAML entry defines the instance name, model, spawn pose
and enabled state. The shared launch rejects duplicate names and derives the
namespace and bridge rules from the same identity. Command-helper profiles are
validated separately because joint order and units depend on the robot
structure. As a result, selection, spawning, commands and feedback remain
associated with the selected instance by construction. The retained live smoke
covers the all-robot profile; automated suites reported in \cref{chap:results}
cover selector rules for the other supported forms.

\section{Configuration-driven robot registry}

Robot instances are selected from YAML entries containing the instance name,
model, spawn pose and enabled flag. The full floor and Room~315 can use
different files while sharing the same loader. The launch layer derives
SDF/URDF paths, ROS namespace and bridge topics from those fields. Users can
select a full name, derived short alias, numeric order, comma-separated set,
\code{all} or \code{none}; the principal surfaces are summarised in
\cref{tab:robots}.

The operator helper uses a separate validated Python profile table for aliases,
joint ordering, command topics and angular or linear unit semantics. Tests
cover both configuration boundaries, although the YAML entries and helper
profiles are not generated from a common source.

\begin{table}[htbp]
\centering
\caption{Principal robot selectors and command surfaces.}
\label{tab:robots}
\begin{tabularx}{\textwidth}{@{}l l l X@{}}
\toprule
\textbf{Selector} & \textbf{Instance} & \textbf{Type} & \textbf{Primary command}\\
\midrule
\code{kuka} & \code{kuka1} & KUKA KR6 &
\makecell[l]{\code{/kuka1/}\\\code{joint\_trajectory}}\\ \tablerowrule
\code{staubli} & \code{staubli1} & Stäubli TX2 &
\makecell[l]{\code{/staubli1/}\\\code{joint\_trajectory}}\\ \tablerowrule
\code{hc10} & \code{yaskawa\_hc10\_1} & Yaskawa HC10 &
\makecell[l]{\code{/yaskawa\_hc10\_1/}\\\code{joint\_trajectory}}\\ \tablerowrule
\code{hc10dt} & \code{yaskawa\_hc10dt\_1} & Yaskawa HC10DT &
\makecell[l]{\code{/yaskawa\_hc10dt\_1/}\\\code{joint\_trajectory}}\\ \tablerowrule
\code{tiago} & \code{tiago1} & Mobile manipulator &
\makecell[l]{\code{/tiago1/}\\\code{joint\_trajectory}\\\code{/tiago1/cmd\_vel}}\\ \tablerowrule
\code{tiago\_base} & \code{tiago\_base1} & Mobile base with torso &
\makecell[l]{\code{/tiago\_base1/}\\\code{joint\_trajectory}\\
\code{/tiago\_base1/cmd\_vel}}\\
\bottomrule
\end{tabularx}
\end{table}

The YAML-driven launch eliminated duplicated launch scripts for every robot
combination. Tests check unique instance names and asset resolution; separate
helper tests check profile selection, joint counts, units and command topics.

\section{Dynamic bridge generation}

Gazebo and ROS~2 use different transports. The launch process derives bridge
configuration for each selected robot and starts \code{ros\_gz\_bridge}
instances. Joint commands travel from ROS~2 to Gazebo; joint state, odometry
and transform information travel back where required. The design follows the
official bridge model, in which YAML connects a ROS topic, Gazebo topic,
the corresponding message types and communication direction
\citep{gazeborosintegration}.

Deriving the bridges from the selected YAML entries keeps renamed namespaces
and bridge configuration synchronised, while per-instance command topics
preserve one-to-one addressing. The launch selection, spawned identity and
bridge topic are therefore all derived from the same YAML entry, as shown in
\cref{fig:robot-flow}.

\begin{figure}[htbp]
\centering
\resizebox{0.98\textwidth}{!}{\begin{tikzpicture}[
  font=\sffamily\normalsize,
  cfg/.style={
    draw=mfjablue,
    very thick,
    rounded corners=1.5mm,
    fill=mfjalight,
    minimum height=13mm,
    text width=42mm,
    align=center,
    font=\sffamily\normalsize
  },
  helper/.style={
    draw=mfjablue,
    rounded corners=1.3mm,
    fill=mfjablue!6,
    minimum height=29mm,
    text width=38mm,
    align=center,
    font=\sffamily\normalsize
  },
  bridge/.style={
    draw=mfjaorange,
    rounded corners=1.3mm,
    fill=mfjaorange!8,
    minimum height=15mm,
    text width=38mm,
    align=center,
    font=\sffamily\normalsize
  },
  sim/.style={
    draw=mfjagreen,
    rounded corners=1.3mm,
    fill=mfjagreen!7,
    minimum height=15mm,
    text width=38mm,
    align=center,
    font=\sffamily\normalsize
  },
  configflow/.style={
    -{Latex[length=2.2mm]},
    very thick,
    draw=mfjablue
  },
  derive/.style={
    -{Latex[length=2.0mm]},
    thick,
    densely dashed,
    draw=mfjaorange
  },
  command/.style={
    -{Latex[length=2.2mm]},
    very thick,
    draw=mfjablue
  },
  feedback/.style={
    -{Latex[length=2.2mm]},
    very thick,
    draw=mfjagreen
  },
  lane/.style={
    draw=mfjablue!55,
    dashed,
    rounded corners=1.8mm,
    fill=mfjalight!28,
    inner xsep=4mm,
    inner ysep=4mm
  },
  lanehead/.style={
    align=center,
    font=\sffamily\normalsize\bfseries,
    text=mfjablue
  }
]
  % One source of truth selects, validates and derives the bridge rules.
  \node[cfg] (registry) {YAML robot registry\\+ selector};
  \node[cfg,right=9mm of registry] (launch)
    {Shared launch\\+ fixed configuration\\checks};
  \node[cfg,right=9mm of launch] (rules)
    {Selected entries\\+ per-instance bridge rules};
  \draw[configflow] (registry) -- (launch);
  \draw[configflow] (launch) -- (rules);

  % Three isolated identities are laid out as vertical command/feedback paths.
  \node[lanehead,below=20mm of registry] (khead)
    {KUKA lane\\isolated identity \texttt{/kuka1}\\[-0.2ex]
     {\color{mfjablue}$\downarrow$ command}\quad
     {\color{mfjagreen}$\uparrow$ feedback}};
  \node[lanehead,below=20mm of launch] (shead)
    {St\"aubli lane\\isolated identity \texttt{/staubli1}\\[-0.2ex]
     {\color{mfjablue}$\downarrow$ command}\quad
     {\color{mfjagreen}$\uparrow$ feedback}};
  \node[lanehead,below=20mm of rules] (thead)
    {TIAGo lane\\isolated identity \texttt{/tiago1}\\[-0.2ex]
     {\color{mfjablue}$\downarrow$ command}\quad
     {\color{mfjagreen}$\uparrow$ feedback}};

  % KUKA: helper output -> bridge -> simulator, with a separate return path.
  \node[helper,below=6mm of khead] (kh) {
    Validated helper output\\[-0.2ex]
    \texttt{/kuka1/}\\[-0.2ex]
    \texttt{joint\_trajectory}};
  \node[bridge,below=13mm of kh] (kb)
    {per-instance\\\texttt{kuka1} bridge};
  \node[sim,below=13mm of kb] (kg)
    {Gazebo\\model/controller\\\texttt{kuka1}};
  \draw[command] ([xshift=7mm]kh.south) -- ([xshift=7mm]kb.north);
  \draw[feedback] ([xshift=-7mm]kb.north) -- ([xshift=-7mm]kh.south);
  \draw[command] ([xshift=7mm]kb.south) -- ([xshift=7mm]kg.north);
  \draw[feedback] ([xshift=-7mm]kg.north) -- ([xshift=-7mm]kb.south);

  % Staubli: the same bounded path, with its own public topic and entity.
  \node[helper,below=6mm of shead] (sh) {
    Validated helper output\\[-0.2ex]
    \texttt{/staubli1/}\\[-0.2ex]
    \texttt{joint\_trajectory}};
  \node[bridge,below=13mm of sh] (sb)
    {per-instance\\\texttt{staubli1} bridge};
  \node[sim,below=13mm of sb] (sg)
    {Gazebo\\model/controller\\\texttt{staubli1}};
  \draw[command] ([xshift=7mm]sh.south) -- ([xshift=7mm]sb.north);
  \draw[feedback] ([xshift=-7mm]sb.north) -- ([xshift=-7mm]sh.south);
  \draw[command] ([xshift=7mm]sb.south) -- ([xshift=7mm]sg.north);
  \draw[feedback] ([xshift=-7mm]sg.north) -- ([xshift=-7mm]sb.south);

  % TIAGo retains separate joint and base command surfaces in one identity.
  \node[helper,below=6mm of thead] (th) {
    Validated joint helper\\[-0.2ex]
    + independent base command\\[-0.2ex]
    \texttt{/tiago1/}\\[-0.2ex]
    \texttt{joint\_trajectory}\\[-0.2ex]
    \texttt{/tiago1/cmd\_vel}};
  \node[bridge,below=13mm of th] (tb)
    {per-instance\\\texttt{tiago1} bridge};
  \node[sim,below=13mm of tb] (tg)
    {Gazebo\\model/controller\\\texttt{tiago1}};
  \draw[command] ([xshift=7mm]th.south) -- ([xshift=7mm]tb.north);
  \draw[feedback] ([xshift=-7mm]tb.north) -- ([xshift=-7mm]th.south);
  \draw[command] ([xshift=7mm]tb.south) -- ([xshift=7mm]tg.north);
  \draw[feedback] ([xshift=-7mm]tg.north) -- ([xshift=-7mm]tb.south);

  % Lane boundaries make it explicit that no global command topic is shared.
  \begin{scope}[on background layer]
    \node[lane,fit=(khead)(kh)(kb)(kg)] (klane) {};
    \node[lane,fit=(shead)(sh)(sb)(sg)] (slane) {};
    \node[lane,fit=(thead)(th)(tb)(tg)] (tlane) {};
  \end{scope}

  % Selected entries derive three independent lanes, not one broadcast path.
  \coordinate (derivebus) at ($(khead.north)!0.5!(thead.north)+(0,7mm)$);
  \draw[derive] (rules.south) -- ++(0,-5mm) -| (derivebus);
  \node[above=1mm of derivebus,fill=white,inner sep=1mm,
    font=\sffamily\normalsize,text=mfjaorange]
    {derive one isolated lane per selected entry};
  \draw[derive] (derivebus) -| (khead.north);
  \draw[derive] (derivebus) -- (shead.north);
  \draw[derive] (derivebus) -| (thead.north);

  \node[below=7mm of sg,align=center,text width=145mm,
    font=\sffamily\normalsize\bfseries,text=mfjablue]
    {Namespace isolation preserves robot identity across command surfaces,
     per-instance bridge, Gazebo entity and feedback.};
\end{tikzpicture}
}
\caption[Configuration-driven multi-robot command and feedback isolation.]
{Configuration-driven multi-robot launch, command and feedback isolation.
Selected YAML entries generate per-instance bridges; the selected
operator-helper profile addresses the corresponding name. There is no shared
global joint-command topic.}
\label{fig:robot-flow}
\end{figure}

\section{Joint-command helper}

Publishing a raw \code{trajectory\_msgs/JointTrajectory} message is cumbersome
and prone to input errors. The \code{robot\_joint\_command.py} operator tool
lists supported names, checks the expected number of positions, maps names to
joints, accepts radians or degrees for angular joints and uses metres for
linear joints. It then publishes only to the selected namespace.

\begin{lstlisting}[language=bash,caption={Examples of the common robot command tool.},label={lst:robot-command}]
ros2 run mfja_robot_control_config robot_joint_command.py --list

ros2 run mfja_robot_control_config robot_joint_command.py \
  kuka --unit deg --positions 34.38 -57.30 63.03 0 34.38 0

ros2 run mfja_robot_control_config robot_joint_command.py \
  staubli --unit rad --positions 0.1 0.4 -0.6 0 0.5 0
\end{lstlisting}

The helper standardises operator input and creates a short trajectory command
with the required joint ordering. Trajectory planning remains separate, as do
collision-free arm motion, synchronisation with the transport cell and
controller-specific operating constraints.

\section{Interface portability}

The public namespaces and typed state boundaries are independent of the current
simulator backend. This separation makes provider substitution possible while
preserving identity, units, timestamps, freshness and controller-feedback
semantics. The implementation and results reported here use the Gazebo providers
identified in \cref{app:reproducibility}.

\section{Results}

The common selector implements \code{none}, one exact name, comma-separated
sets and \code{all}. In the retained all-robot cold-start smoke, the full-floor
entry point spawned the six enabled registry entries. Each exposed a distinct
joint command topic and a fresh namespaced state sample; both TIAGo variants
also exposed their independent base-velocity topic. The smoke recorded startup
and feedback without issuing motion commands. The command helper and launch-selection rules are covered by the
named automated suites in \cref{chap:results}. The resulting access layer
therefore brings the separate simulation models under a common, maintainable
interface. The live check covers the \code{all} selector, while the automated
tests exercise the remaining selector rules.
